\documentclass[11pt]{article}
\usepackage{preamble}
\renewcommand{\answer}[1]{}

\begin{document}

\ladderheading{AP Statistics \textperiodcentered\ Unit 1 \textperiodcentered\ Topic 1.12 \textperiodcentered\ B: 2026-10-01 \textperiodcentered\ E: 2026-10-09}{Two Surveys of Park Fees}

\focusbanner{TODAY'S PROBLEM}

Suppose a town studies support for a park fee.\par\smallskip \textbf{Mail survey:} A neutral question is mailed to 500 randomly chosen households. Only 120 return it; $78\%$ of respondents oppose the fee.\par \textbf{Phone survey:} The town calls 300 random phone numbers. Of 240 people who answer, $62\%$ oppose after hearing, \emph{Do you support a NEW FEE that would tax park users?}\par\smallskip One council member says $62\%$ must be closer to the truth because more people answered. Another says $78\%$ must be closer because the households were randomly chosen.\par
\begin{center}
\textbf{Park-fee surveys}\par\smallskip
\begin{tabular}{|l|c|c|c|}
\hline
\textbf{} & \textbf{Selected} & \textbf{Replied} & \textbf{Oppose} \\ \hline
\textbf{Mail} & 500 & 120 & $78\%$ \\ \hline
\textbf{Phone} & 300 & 240 & $62\%$ \\ \hline
\end{tabular}
\end{center}
\begin{firsttakebox}{FIRST TAKE --- before the video (5 min)}
Which estimate, if either, would you report to the town? Name one detail behind your choice.
\workspace{0.9in}
\end{firsttakebox}

\begin{callout}{\IconBook\ BIAS FAVORS SOME RESPONSES (keep on the page)}{calloutgreen}
\textbf{Bias} systematically favors some responses over others. \textbf{Nonresponse bias} is possible when
selected people who do not respond differ from those who do; the response rate is
$\text{responses}/\text{selected}$. \textbf{Response bias} can come from confusing or leading wording.
A larger sample or response rate does not repair a different design flaw. Match the fix to the source of bias.
\end{callout}

\newpage
\textbf{Finish after the video:}\par\smallskip

\dokbadge{1}~\textbf{(a)} Find each survey's response rate.\par
\workspace{0.8in}

\dokbadge{2}~\textbf{(b)} In two or three sentences, explain one possible bias in each survey. Say which direction each could push opposition, or why its direction is unknown.\par
\workspace{1.4in}

\dokbadge{3}~\textbf{(c)} Has either council member shown which estimate is closer to the town's true opinion? Use one detail from each survey to justify your answer. Choose one survey to improve, and explain how your change addresses its problem.\par
\workspace{2.0in}

\begin{sentenceframebox}\raggedright\textbf{Frame:} I cannot tell which is closer because \blank[1.7in]. In the mail survey \blank[1.8in]; in the phone survey \blank[1.8in].\end{sentenceframebox}

\begin{sentenceframebox}\raggedright\textbf{Frame:} I would change \blank[1.8in] because that addresses \blank[1.8in].\end{sentenceframebox}

\begin{summaryexitbox}{EXIT --- turn this in}
One thing the video changed about my first take: \hrulefill\par
\workspace{0.4in}
\end{summaryexitbox}

\end{document}
